% Standalone section; requires amsmath, amssymb, amsthm, hyperref.
% Uses the theorem, proposition, lemma, proof, remark environments.
% No custom mathematical macros or external labels are required.
\section{Exact finite-dimensional orbit-span optimization}
\label{sec:exact-finite-rank}

The article's question \texttt{spn:q:selection} asks, in part, for the
minimum total orbit-span rank at a fixed quadruple of parameters
$(d,s,u,v)$. Its theorem \texttt{spn:thm:greedy} gives the sharp dimension
fraction $(1-u)^2/(1+v)$, while its sharpness construction varies the
dimensions and takes $s=1$. The result below determines the entire set of
possible total ranks with both dimensions fixed. It is an unconditional
finite-dimensional theorem extracted from the spanning interface of the
GRH development; it makes no assumption about zeros of $L$-functions.

Let $S,T$ be finite-dimensional complex Hilbert spaces, let a finite group
$G$ act unitarily on $T$ through $\rho$, and let $W\subseteq T$ be
$G$-invariant. Write $s=\dim S>0$, $d=\dim W>0$, and let $P$ be the
orthogonal projection onto $W$. For $B:S\to T$, assume
\begin{align}
 \|PBP_U\|_{\mathrm{HS}}^2&\ge \dim U-us
       &&\text{for every subspace }U\subseteq S,\label{rank:eq:near}\\
 \frac{d}{|G|}\sum_{g\in G}\|B^*\rho(g)B\|_{\mathrm{HS}}^2
       &\le(1+v)s^2.\label{rank:eq:moment}
\end{align}
Here $P_U$ is the orthogonal projection onto $U$, the Hilbert--Schmidt
norm is the square root of the sum of the squared singular values,
and the parameter domain is exactly
\[
             0\le u<1,\qquad v\ge0.
\]
The object to minimize is
\[
 F=\sum_{g\in G}\operatorname{im}(P\rho(g)B),\qquad r=\dim F.
\]
The spaces, group, representation and map may vary, while $d,s,u,v$
are fixed. Put $a=1-u$ and $z=1+v$.

\begin{theorem}[The complete feasible rank spectrum]
\label{rank:thm:spectrum}
There exists a system satisfying \eqref{rank:eq:near}--\eqref{rank:eq:moment}
if and only if $as\le d$. In the feasible case the set of all possible
integers $\dim F$ is precisely
\[
 \{r_*,r_*+1,\ldots,d\},\qquad
 r_* =\max\left\{\left\lceil as\right\rceil,
                    \left\lceil\frac{a^2d}{z}\right\rceil\right\}.
\]
Every allowed rank $r$ can be attained with $T=W=\mathbb C^d$, $P=I$,
a cyclic group of order $r$, a contraction $B$, and
$\operatorname{rank}B=\lceil as\rceil$. The construction works for
arbitrary real $u,v$ in the displayed parameter domain.
\end{theorem}

\begin{proof}
First establish the necessary conditions. Applying
\eqref{rank:eq:near} to $\ker(PB)$ gives
\[
 0\ge s-\operatorname{rank}(PB)-us,
 \qquad r\ge\operatorname{rank}(PB)\ge as.
\]
In particular $as\le d$ is necessary. Define the positive semidefinite
operator
\[
 M=\frac1{|G|}\sum_{g\in G}\rho(g)BB^*\rho(g)^*\quad\text{on }T,
 \qquad A=(PMP)|_W.
\]
Since $P$ commutes with the representation, $A$ is the covariance of
the maps $P\rho(g)B$, has image $F$, and
\[
 \operatorname{tr}A=\|PB\|_{\mathrm{HS}}^2\ge as.
\]
Indeed a vector belongs to $\ker A$ exactly when it is perpendicular
to every image $\operatorname{im}(P\rho(g)B)$. Averaging a double sum,
using unitarity, and then substituting $g^{-1}h$ yields the identity
\[
 \operatorname{tr}(M^2)
   =\frac1{|G|}\sum_{g\in G}\|B^*\rho(g)B\|_{\mathrm{HS}}^2.
\]
Relative to $T=W\oplus W^\perp$, write
$M=\left(\begin{smallmatrix}A&C\\C^*&D\end{smallmatrix}\right)$.
Consequently
\[
 \operatorname{tr}(M^2)=\operatorname{tr}(A^2)
              +2\|C\|_{\mathrm{HS}}^2+\operatorname{tr}(D^2)
             \ge\operatorname{tr}(A^2).
\]
The ordinary Cauchy--Schwarz inequality for the $r$ positive
eigenvalues of $A$ therefore gives
\[
 (as)^2\le(\operatorname{tr}A)^2
       \le r\operatorname{tr}(A^2)
       \le r\operatorname{tr}(M^2)
       \le \frac{rzs^2}{d}.
\]
Thus $r\ge a^2d/z$ as well, proving both lower bounds.

We now construct a system for every integer $r$ satisfying
$r_*\le r\le d$. Put $k=\lceil as\rceil$. Since $0<a\le1$,
\[
              1\le k\le\min(s,r).
\]
Choose $\lambda_1,\ldots,\lambda_k\in(0,1]$ with
$\sum_j\lambda_j=as$: for instance take the first $k-1$ equal to one
and the last equal to $as-k+1$. Let $E:\mathbb C^s\to\mathbb C^r$
send $e_j$ to $\sqrt{\lambda_j}\,e_j$ for $j\le k$, and to zero for
$j>k$. Let $Q$ be the unitary Fourier matrix of order $r$,
\[
 Q_{ij}=r^{-1/2}\exp\bigl(2\pi\mathrm{i}(i-1)(j-1)/r\bigr),
          \qquad 1\le i,j\le r,
\]
and set $B=QE$, with the result embedded in the first $r$ coordinates
of $\mathbb C^d$. We take $T=W=\mathbb C^d$ and $P=I$.
Because $B^*B\le I_S$ and $\operatorname{tr}(B^*B)=as$,
\[
 \operatorname{tr}(I_S-B^*B)=us.
\]
For every $U\subseteq S$, positivity implies
\[
 \|BP_U\|_{\mathrm{HS}}^2
 =\dim U-\operatorname{tr}\bigl(P_U(I_S-B^*B)\bigr)
 \ge\dim U-us.
\]
This checks the full subspace condition, not only its instance $U=S$.

Let $G=\mathbb Z/r\mathbb Z$, put $\omega=\exp(2\pi\mathrm{i}/r)$,
and let the generator act as
$\operatorname{diag}(1,\omega,\ldots,\omega^{r-1})$ on the first
$r$ coordinates and as the identity on the others. Distinct Fourier
characters make every off-diagonal entry of the averaged covariance
vanish. Every diagonal entry on the first $r$ coordinates is $as/r$,
because every entry of $Q$ has absolute value $r^{-1/2}$. Thus
\[
 \frac1{|G|}\sum_{g\in G}\rho(g)BB^*\rho(g)^*
       =\frac{as}{r}\,P_{\mathbb C^r}.
\]
This operator has rank $r$, so the orbit images span exactly that
$r$-dimensional coordinate subspace. The covariance identity gives
\[
 \frac d{|G|}\sum_{g\in G}\|B^*\rho(g)B\|_{\mathrm{HS}}^2
       =\frac{d(as)^2}{r}\le zs^2,
\]
as required. When $r=1$, the same construction uses the trivial group
and the $1\times1$ Fourier matrix. Finally, $a^2d/z\le d$, so
$as\le d$ is also sufficient for $r_*\le d$.
\end{proof}

\begin{proposition}[Rigidity when the second-moment bound is exact]
\label{rank:prop:equality}
Suppose a feasible system attains
$r=a^2d/z$ exactly. Then, and only then, its map and orbit covariance
have all three properties
\[
 \operatorname{im}B\subseteq W,\qquad
 0\le B^*B\le I_S,\quad\operatorname{tr}(B^*B)=as,
 \qquad M=\frac{as}{r}P_F.
\]
In the converse direction $F$ is the orbit span and is assumed to have
dimension $r=a^2d/z$. In particular an equality example can exist only
when this number is an integer at least $\lceil as\rceil$; when those
conditions hold, Theorem~\ref{rank:thm:spectrum} supplies one.
\end{proposition}

\begin{proof}
All inequalities in the lower-bound chain must be equalities.
Equality in the covariance compression implies $C=D=0$ in the block
decomposition of $M$. Since $W$ is invariant,
\[
 0=\operatorname{tr}D
  =\frac1{|G|}\sum_g\|(I-P)\rho(g)B\|_{\mathrm{HS}}^2
  =\|(I-P)B\|_{\mathrm{HS}}^2,
\]
so $B$ takes its values in $W$. Equality in Cauchy--Schwarz makes all
positive eigenvalues of $A$ equal, and equality in the trace bound
makes their sum $as$. Hence $M=(as/r)P_F$ and
$\operatorname{tr}(B^*B)=as$.

It remains to verify the contraction assertion. Let $\mu_1,\ldots,\mu_s$
be the eigenvalues of $B^*B$, and apply \eqref{rank:eq:near} to the
span of the eigenvectors with $\mu_j<1$. This yields
\[
 \sum_j(1-\mu_j)_+\le us.
\]
But $\sum_j(1-\mu_j)=us$. Subtracting these identities shows
$\sum_j(\mu_j-1)_+=0$, so every $\mu_j\le1$.
Conversely, a contraction of trace $as$ satisfies
\eqref{rank:eq:near} by the positivity argument in the construction,
and the stated covariance gives equality in
\eqref{rank:eq:moment}.
\end{proof}

\paragraph{What this changes.}
There are two independent obstructions: a single source map already
needs at least $\lceil(1-u)s\rceil$ independent image directions, while
the second moment forces at least
$\lceil(1-u)^2d/(1+v)\rceil$ orbit directions. The larger obstruction is
always attainable, and every larger rank is attainable too. For
$(d,s,u,v)=(100,90,1/10,1/5)$ the two bounds are $81$ and $68$,
respectively, and the exact answer is $81$. For
$(100,10,1/10,1/5)$ they are $9$ and $68$, and the answer is $68$.
For $(100,120,1/10,1/5)$ no admissible system exists, because
$(1-u)s=108>d$.

The covariance lower bound is classical effective-rank mathematics;
Fourier equalization is a standard tight-frame construction. The
contribution here is their exact combination for the article's
fixed-parameter extremal problem, including feasibility, the complete
integer spectrum, and equality conditions. This resolves the
total-rank portion of \texttt{spn:q:selection}. It does not determine
the optimal number of translates needed to reach a prescribed
intermediate rank, which is the other part of that question.
